% Conformité guide ENSIASD (template rapport de stage).
% Compilation : xelatex -interaction=nonstopmode main.tex  (deux fois)

\documentclass[12pt,a4paper,oneside]{report}

\usepackage{fontspec}
\usepackage{polyglossia}
\setmainlanguage{french}
\setmainfont{Times New Roman}[
  Ligatures=TeX,
  SmallCapsFont={Times New Roman},
  SmallCapsFeatures={Letters=Uppercase,Scale=0.82}
]

% Guillemets style babel-french (\og \fg), absents de polyglossia.
\providecommand{\og}{\guillemotleft~}
\providecommand{\fg}{\unskip~\guillemotright}

\setlength{\emergencystretch}{3em}
\tolerance=2000

\usepackage[a4paper,top=2.5cm,bottom=2.5cm,left=2.5cm,right=2.5cm,headheight=18pt,headsep=8pt,footskip=22pt]{geometry}
\usepackage{setspace}
\onehalfspacing

\usepackage{amsmath,amssymb}
\usepackage{graphicx}
\usepackage{needspace}
\usepackage{array,booktabs,tabularx,longtable}
\usepackage{xcolor}
\usepackage{colortbl}
\definecolor{ensiBlue}{RGB}{0, 90, 156}
\definecolor{tableHeadBlue}{RGB}{0, 90, 156}
\newcommand{\TableHead}{\rowcolor{tableHeadBlue}}
\newcommand{\tht}[1]{\textcolor{white}{\textbf{#1}}}
% Bloc techno : texte à gauche, logo à droite, espacement élastique
% pour remplir la page sans trou en bas.
\newcommand{\TechBloc}[3]{%
  \par
  \vspace{0.32em plus 0.2em minus 0.08em}%
  {\color{ensiBlue!35}\hrule height 0.45pt}%
  \vspace{0.28em plus 0.12em}%
  \noindent
  \begin{tabular}{@{}p{\dimexpr\textwidth-1.75cm} @{\hspace{0.28cm}} c @{}}
    \textbf{#1.} #3 &
    \raisebox{-0.42\height}{\includegraphics[height=1.05cm]{#2}}
  \end{tabular}%
  \par
}
\usepackage{microtype}
\usepackage{csquotes}
\usepackage{enumitem}
\setlist{nosep,leftmargin=1.5em}

\usepackage[explicit]{titlesec}
\usepackage{fancyhdr}
\usepackage{caption}
\usepackage{tocloft}

% En-tête ENSIASD : libellé italique gras à gauche + filet.
% Chaque style redéfinit entièrement têtes et pieds (fancyhdr).
\newcommand{\EnsiHeadText}{\small\bfseries\itshape\nouppercase{\leftmark}}
\newcommand{\EnsiPageStyle}{%
  \fancyhf{}%
  \renewcommand{\headrulewidth}{0.45pt}%
  \renewcommand{\footrulewidth}{0pt}%
  \fancyhead[L]{\EnsiHeadText}%
  \fancyfoot[R]{\thepage}%
}
\fancypagestyle{plain}{\EnsiPageStyle}
\fancypagestyle{frontmatter}{\EnsiPageStyle}
\fancypagestyle{mainmatter}{\EnsiPageStyle}
\pagestyle{mainmatter}

\renewcommand{\chaptermark}[1]{}
\renewcommand{\sectionmark}[1]{}

\renewcommand{\chaptername}{CHAPITRE}
\newcommand{\EnsiChapterTitle}[1]{%
  \markboth{Chapitre \thechapter{} : #1}{}%
  \MakeUppercase{\chaptername~\thechapter{} : #1}%
}
\newcommand{\EnsiChapterStar}[1]{%
  \markboth{#1}{}%
  \MakeUppercase{#1}%
}
\titleformat{\chapter}[block]
  {\normalfont\bfseries\fontsize{16}{20}\selectfont\raggedleft}
  {}
  {0pt}
  {\EnsiChapterTitle{#1}}
\titlespacing*{\chapter}{0pt}{0pt}{54pt}

\titleformat{name=\chapter,numberless}[block]
  {\normalfont\bfseries\fontsize{16}{20}\selectfont\raggedleft}
  {}
  {0pt}
  {\EnsiChapterStar{#1}}
\titlespacing*{name=\chapter,numberless}{0pt}{0pt}{54pt}

\renewcommand{\thesection}{\Roman{section}}
\titleformat{\section}
  {\normalfont\bfseries\fontsize{14}{18}\selectfont}
  {\thesection.}
  {0.6em}
  {#1}
\titlespacing*{\section}{0pt}{12pt}{12pt}

\renewcommand{\thesubsection}{\thesection.\arabic{subsection}}
\titleformat{\subsection}
  {\normalfont\bfseries\fontsize{12}{16}\selectfont}
  {\thesubsection}
  {0.5em}
  {#1}
\titlespacing*{\subsection}{0pt}{6pt}{6pt}

\renewcommand{\thesubsubsection}{\thesection.\arabic{subsection}.\arabic{subsubsection}}
\titleformat{\subsubsection}
  {\normalfont\bfseries\fontsize{12}{16}\selectfont}
  {\thesubsubsection}
  {0.5em}
  {#1}
\titlespacing*{\subsubsection}{0pt}{6pt}{6pt}

\usepackage{newfloat}
\DeclareFloatingEnvironment[
  fileext=los,
  listname={Liste des schémas},
  name=Schéma,
  placement=htbp
]{schema}

\renewcommand{\thefigure}{\arabic{chapter}.\arabic{figure}}
\renewcommand{\thetable}{\arabic{chapter}.\arabic{table}}
\renewcommand{\theschema}{\arabic{chapter}.\arabic{schema}}

\DeclareCaptionLabelFormat{ensi}{#1 : #2}
\captionsetup[figure]{
  name={Figure},
  font={it},
  justification=centering,
  skip=6pt,
  labelformat=ensi,
  labelsep=colon
}
\captionsetup[schema]{
  font={it},
  justification=centering,
  skip=6pt,
  labelformat=ensi,
  labelsep=colon
}
% Tableaux : titre en haut, 11 pt, centré
\captionsetup[table]{
  font={it},
  justification=centering,
  skip=6pt,
  labelformat=ensi,
  labelsep=colon,
  position=top
}
\usepackage{etoolbox}
\AtBeginEnvironment{figure}{\fontsize{11}{13}\selectfont}
\AtBeginEnvironment{table}{\fontsize{11}{13}\selectfont}
\AtBeginEnvironment{schema}{\fontsize{11}{13}\selectfont}
\AfterEndPreamble{%
  \renewcommand{\figurename}{Figure}%
  \renewcommand{\tablename}{Tableau}%
}

\graphicspath{{figures/}{../rapport/}}

\newcommand{\FigPlaceholder}[1]{%
  \begin{center}
    \fbox{\parbox{0.88\textwidth}{\centering\vspace{2.2cm}\textit{#1}\vspace{2.2cm}}}
  \end{center}
}

\setcounter{tocdepth}{3}
\setcounter{secnumdepth}{3}

\setlength{\parindent}{1.25em}
\setlength{\parskip}{0pt}

\usepackage{xurl}
\usepackage{hyperref}
\usepackage{bookmark}

\AtBeginDocument{%
  \renewcommand{\contentsname}{Table des matières}%
  \renewcommand{\listfigurename}{Liste des figures}%
  \renewcommand{\listtablename}{Liste des tableaux}%
  \renewcommand{\bibname}{Références bibliographiques}%
  \renewcommand{\figurename}{Figure}%
  \renewcommand{\tablename}{Tableau}%
}

\makeatletter
\renewcommand\tableofcontents{%
  \chapter*{\contentsname}%
  \@starttoc{toc}%
}
\renewcommand\listoffigures{%
  \chapter*{\listfigurename}%
  \@starttoc{lof}%
}
\renewcommand\listoftables{%
  \chapter*{\listtablename}%
  \@starttoc{lot}%
}
\makeatother
